\documentclass[11pt,a4paper]{amsart}
\usepackage[T1]{fontenc}
\usepackage{lmodern}
\usepackage{amssymb}
\usepackage{microtype}
\usepackage{booktabs}
\usepackage{graphicx}
\usepackage{tikz}
\usepackage[hyphens]{url}
\usepackage[colorlinks=true,linkcolor=black,citecolor=black,urlcolor=black]{hyperref}
\usepackage[margin=2.5cm]{geometry}

\hypersetup{
  pdftitle={A new upper bound for the density of planar sets avoiding unit distances},
  pdfauthor={Sergi Gal\'an},
  pdfsubject={Planar sets avoiding unit distances},
  pdfkeywords={unit distance graph, distance avoiding sets, density, fractional chromatic number, chromatic number of
  the plane, linear programming, inclusion-exclusion}}

\newtheorem{theorem}{Theorem}
\newtheorem{corollary}[theorem]{Corollary}
\newtheorem{lemma}[theorem]{Lemma}
\newtheorem{proposition}[theorem]{Proposition}
\theoremstyle{remark}
\newtheorem*{remark}{Remark}

\newcommand{\Q}{\mathbb{Q}}
\newcommand{\R}{\mathbb{R}}
\newcommand{\Z}{\mathbb{Z}}
\newcommand{\E}{\mathbb{E}}
\newcommand{\chif}{\chi_{\mathrm f}}
% the numbers of the main certificate (filled in from the certificate file)
\newcommand{\NPTS}{30}
\newcommand{\NNEW}{7}
\newcommand{\MONE}{0.24380}
\newcommand{\INVM}{4.1017}
\newcommand{\NUNIT}{80}
\newcommand{\NATOMS}{84\,893}
\newcommand{\NCONG}{15\,709}

\title{A new upper bound for the density of planar sets avoiding unit distances}
\author{Sergi Gal\'an}
\thanks{These results were obtained with the assistance of an AI system (Claude, by Anthropic), which found the
configuration, wrote the programs, ran the checks and helped to draft this note; a second program and a referee report
were produced by separate AI agents. The method and the 23-point configuration are those of Ambrus, Csisz\'arik,
Matolcsi, Varga and Zs\'amboki \cite{ACMVZ}. No mathematician has checked the note yet. The author is not a
professional mathematician and submits it for independent verification.}
\date{Draft, 11 October 2026}
\subjclass[2020]{Primary 52C10, 05C15; Secondary 42B10, 90C05}
\keywords{unit distance graph, distance avoiding sets, density, fractional chromatic number, chromatic number of the
plane, linear programming bounds, inclusion--exclusion}

\begin{document}

\begin{abstract}
Let $m_1(\R^2)$ be the supremum of the upper densities of measurable planar sets in which no two points are at
distance $1$. Ambrus, Csisz\'arik, Matolcsi, Varga and Zs\'amboki proved $m_1(\R^2)\le0.2470$, confirming Erd\H{o}s's
conjecture $m_1(\R^2)<1/4$, with a linear program whose constraints come from the patterns in which a random copy of a
set $X_{23}$ of $23$ points meets the set; their computer search for $X$ stopped at $24$ points. Adding $\NNEW$ points
to $X_{23}$, with coordinates in the same field $\Q(\sqrt3,\sqrt{11})$, we prove $m_1(\R^2)\le\MONE$. This improves the published bound; the same group has announced $0.2415$ without a
published proof. Our proof is a rational certificate for the program, checked in exact and interval arithmetic by two
programs written separately.
With the transfer theorem of OpenAI, as extended to fractional colourings in an earlier note, the fractional chromatic
number of the plane is therefore at least $\INVM$. In floating point the value of the program keeps decreasing as points
are added, so the method has not reached its limit.
\end{abstract}

\maketitle

\section{The result}

A set $A\subseteq\R^2$ \emph{avoids the distance $1$} if $|x-y|\ne1$ for all $x,y\in A$. Its upper density is
$\overline\delta(A)=\limsup_{R\to\infty}|A\cap B(0,R)|/|B(0,R)|$, and
\[
m_1(\R^2)=\sup\{\overline\delta(A):A\subseteq\R^2\text{ measurable, avoiding the distance }1\}.
\]
The question was popularized by Leo Moser; Croft's construction gives $m_1(\R^2)\ge0.22936$ \cite{Croft}, and no better
construction is known. Erd\H{o}s wrote that ``it seems very likely that $m_1(\R^2)$ is less than $1/4$''
\cite{Erdos85}. The upper bounds came from unit-distance graphs with large fractional chromatic number (the Moser spindle
gives $2/7$; Bellitto, P\^echer and S\'edillot $0.2565$ \cite{BPS}) and from linear programs on the autocorrelation of
$A$: Sz\'ekely $12/43\approx0.2791$ \cite{Szekely}, Oliveira Filho and Vallentin $0.2684$ \cite{OFV}, Keleti, Matolcsi,
Oliveira Filho and Ruzsa $0.2588$ \cite{KMOR}, Ambrus and Matolcsi $0.2544$ \cite{AM}, and finally Ambrus,
Csisz\'arik, Matolcsi, Varga and Zs\'amboki, $m_1(\R^2)\le0.2470$ \cite{ACMVZ}, which settled Erd\H{o}s's conjecture.
In December 2023 Varga announced that the group had improved this to $0.2415$ \cite{Varga}; we found no paper, preprint
or certificate for that bound.
Their program combines both approaches: for a finite set $X$, the probabilities of the patterns in which a random
congruent copy of $X$ meets $A$ satisfy linear equations (``inclusion--exclusion''), whose coefficients are values of
the autocorrelation of $A$. They found a set $X_{23}$ of $23$ points by a beam search over sets of at most $24$ points
(a week on $128$ processors \cite[§9]{ACMVZ}).

\begin{theorem}\label{thm:main}
Every measurable planar set avoiding the distance $1$ has upper density at most $\MONE$. That is,
$m_1(\R^2)\le\MONE$.
\end{theorem}

This is weaker than the announced $0.2415$, but, as far as we know, it is the best upper bound with a published proof.
The method is entirely that of \cite{ACMVZ}; what is new is the configuration and the certificate. The proof is the
program of \cite{ACMVZ} for a set $X_{\NPTS}$ of $\NPTS$ points: $X_{23}$ and $\NNEW$ more points in
the same field $\Q(\sqrt3,\sqrt{11})$ (Section~\ref{sec:config}). Section~\ref{sec:ie} states the certificate that
proves the bound and Section~\ref{sec:checks} how it was checked. Section~\ref{sec:search} describes how the points were
found and what larger sets give in floating point.

\subsection*{Consequences}
Write $\chif(\R^2)$ for the fractional chromatic number of the unit-distance graph of the plane, that is, the supremum of
$\chif(G)$ over finite unit-distance graphs $G$. Averaging translates of a set avoiding $1$ over the vertices of $G$
gives $m_1(\R^2)\le1/\chif(G)$; the converse direction is the hard one.

\begin{corollary}\label{cor:frac}
$\chif(\R^2)\ge1/m_1(\R^2)\ge\INVM$, and the measurable fractional chromatic number of the plane is at least $\INVM$.
For every $\varepsilon>0$ and every large $n$ there are $n$-point unit-distance graphs whose independence number is
below $(\MONE+\varepsilon)n$.
\end{corollary}

\begin{proof}
The measurable statement is the averaging argument above in the measurable setting. That $\chif(\R^2)=1/m_1(\R^2)$, and
that the least independence ratio of $n$-point unit-distance graphs tends to $m_1(\R^2)$, is proved in
\cite[Cors.~2 and~3]{Galan} from the transfer theorem of OpenAI \cite[Thm.~2.3]{OpenAI}; insert
Theorem~\ref{thm:main}.
\end{proof}

Before, these bounds were $4.0485$ and $0.2470$ \cite{ACMVZ,Galan}; Dúcz and Varga had found a finite unit-distance
graph with fractional chromatic number above $4.0007$ \cite{DV}. The same substitution improves the other bounds of
\cite{Galan} that use $m_1(\R^2)$: for instance, for suitable sets $D$ of $k$ integers, $\chi(\R^2,D)\ge
m_1(\R^2)^{-k}>\INVM^k$ \cite[Cor.~6]{Galan}.

\section{The inclusion--exclusion bound}\label{sec:ie}

Let $X=\{x_1,\dots,x_n\}\subseteq\R^2$, $r_{ij}=|x_i-x_j|$, let $E_1$ be the set of pairs $\{i,j\}$ with $r_{ij}=1$, and
let $\mathcal S$ be the family of subsets of $\{1,\dots,n\}$ containing no pair of $E_1$ (the independent sets of the
unit-distance graph of $X$, including $\emptyset$). Let $\mathcal C$ be a list of pairs $(I,J)$ of sequences of distinct
indices of the same length with $|x_{I_a}-x_{I_b}|=|x_{J_a}-x_{J_b}|$ for all $a,b$.

\begin{proposition}[Ambrus, Csisz\'arik, Matolcsi, Varga, Zs\'amboki \cite{ACMVZ}]\label{prop:ie}
Suppose that real numbers $\psi_0,\psi_T,\lambda_i$ $(1\le i\le n)$, $u_{ij}$ $(\{i,j\}\notin E_1)$ and $w_c$
$(c\in\mathcal C)$ satisfy
\begin{gather*}
\Phi(t)=\psi_0J_0(t)-\sum_i\lambda_i-\sum_{i<j}u_{ij}J_0(r_{ij}t)\ \ge\ 1\qquad(t\ge0),\\
Q(S)=\psi_T+\sum_{i\in S}\lambda_i+\sum_{i<j,\ i,j\in S}u_{ij}+\sum_{(I,J)\in\mathcal C}w_{(I,J)}\bigl([I\subseteq S]-[J\subseteq S]\bigr)
\ \ge\ 0\qquad(S\in\mathcal S),
\end{gather*}
where the sums of $u_{ij}$ run over the pairs $i<j$ with $\{i,j\}\notin E_1$. Then $m_1(\R^2)\le\psi_T$.
\end{proposition}

\begin{proof}[Proof (after \cite{ACMVZ})]
It suffices to bound the density $\delta$ of a measurable set $A$ avoiding $1$ that is periodic for a lattice
$\Lambda$ \cite[§2]{ACMVZ}. Let $c_\xi$ ($\xi\in\Lambda^*$) be the Fourier coefficients of $1_A$ on $\R^2/\Lambda$, and
$\kappa$ the image of the measure $\sum_\xi|c_\xi|^2\delta_\xi$ under $\xi\mapsto2\pi|\xi|$: a measure on $[0,\infty)$ of
total mass $\delta$ (Parseval) with $\kappa(\{0\})=\delta^2$. The radial average of the autocorrelation
$\varphi(v)=\delta(A\cap(A-v))$ is $f(r)=\int J_0(rt)\,d\kappa(t)$, and $f(1)=0$ since $A$ avoids $1$. Choose $y$
uniformly in $\R^2/\Lambda$ and $g\in O(2)$ according to Haar measure, and let $S=\{i:y+gx_i\in A\}$. Then
$S\in\mathcal S$, $\Pr(i\in S)=\delta$, $\Pr(i,j\in S)=f(r_{ij})$, and $\Pr(I\subseteq S)=\Pr(J\subseteq S)$ for
$(I,J)\in\mathcal C$: some isometry $h=h_0+c$ maps $x_{I_a}$ to $x_{J_a}$ for every $a$, and $(y,g)\mapsto(y+gc,gh_0)$
preserves the distribution of $(y,g)$ (reflections are allowed because $g$ ranges over $O(2)$; $9\,474$ of the $\NCONG$
relations of our certificate for $X_{\NPTS}$ need an isometry that reverses orientation). Hence
$0\le\E\,Q(S)=\psi_T+\delta\sum_i\lambda_i+\sum_{i<j}u_{ij}f(r_{ij})$. Integrating $\Phi\ge1$ against $\kappa$,
including the atom at $0$,
\[
\delta=\int d\kappa\le\int\Phi\,d\kappa=\psi_0f(1)-\delta\sum_i\lambda_i-\sum_{i<j}u_{ij}f(r_{ij})\le\psi_T.\qedhere
\]
\end{proof}

The best such bound for given $X$ and $\mathcal C$ is the value of a linear program, whose dual variables are the
probabilities of the patterns and the measure $\kappa$; for $X_{23}$ it is $0.2469726\ldots$ \cite[§8]{ACMVZ}, and we
compute $0.24697261$ for the same program, discretized as in \cite{ACMVZ}.

\section{The configuration}\label{sec:config}

Write points as complex numbers. $X_{\NPTS}$ consists of the $23$ points of \cite[Table~1]{ACMVZ}, in their order, and the
$\NNEW$ points of Table~\ref{tab:new}. All coordinates lie in $\Q(\sqrt3,\sqrt{11})$; we compute in this field exactly,
with the basis $1,\sqrt3,\sqrt{11},\sqrt{33}$. The unit-distance graph of $X_{\NPTS}$ has $\NUNIT$ edges (that of $X_{23}$
has $47$) and $\NATOMS$ independent sets.

\begin{table}[ht]
\centering
\renewcommand{\arraystretch}{1.35}
\begin{tabular}{rlll}
\toprule
$k$ & $\operatorname{Re}x_k$ & $\operatorname{Im}x_k$ & points at distance $1$\\
\midrule
$24$ & $\frac{3-\sqrt{33}}{6}$ & $\frac{\sqrt3}{3}$ & 3, 15, 20, 22, 25, 27 \\
$25$ & $\frac{11-\sqrt{33}}{12}$ & $\frac{-\sqrt3+\sqrt{11}}{12}$ & 6, 15, 21, 22, 24, 30 \\
$26$ & $\frac{1}{6}$ & $-\frac{\sqrt{11}}{6}$ & 2, 11, 12, 21 \\
$27$ & $\frac{9-\sqrt{33}}{6}$ & $\frac{\sqrt3}{3}$ & 4, 7, 10, 16, 20, 24 \\
$28$ & $\frac{4}{3}$ & $\frac{3\sqrt3+\sqrt{11}}{6}$ & 3, 5, 17, 18, 22 \\
$29$ & $\frac{2-\sqrt{33}}{6}$ & $\frac{2\sqrt3+\sqrt{11}}{6}$ & 12, 13, 18, 19, 22 \\
$30$ & $\frac{17-\sqrt{33}}{12}$ & $\frac{5\sqrt3+\sqrt{11}}{12}$ & 2, 6, 12, 25 \\
\bottomrule
\end{tabular}
\caption{The points $x_{24},\dots,x_{30}$ added to $X_{23}$, and the points of $X_{30}$ at distance $1$ from each.}\label{tab:new}
\end{table}

\begin{figure}[ht]
\centering
\begin{tikzpicture}[scale=4.6,every node/.style={font=\scriptsize}]
\draw[black!35,line width=0.4pt] (0.00000,0.00000) -- (1.00000,0.00000);
\draw[black!35,line width=0.4pt] (0.00000,0.00000) -- (0.50000,0.86603);
\draw[black!35,line width=0.4pt] (0.00000,0.00000) -- (0.83333,0.55277);
\draw[black!35,line width=0.4pt] (0.00000,0.00000) -- (-0.06205,0.99807);
\draw[black!35,line width=0.4pt] (0.00000,0.00000) -- (0.72871,-0.68482);
\draw[black!35,line width=0.4pt] (1.00000,0.00000) -- (0.50000,0.86603);
\draw[black!35,line width=0.4pt] (1.00000,0.00000) -- (1.50000,0.86603);
\draw[black!35,line width=0.4pt] (1.00000,0.00000) -- (0.10462,0.44530);
\draw[black!35,line width=0.4pt] (1.00000,0.00000) -- (0.27129,0.68482);
\draw[black!35,line width=0.4pt] (1.00000,0.00000) -- (0.04257,-0.28868);
\draw[red!70!black,line width=0.5pt] (1.00000,0.00000) -- (0.16667,-0.55277);
\draw[red!70!black,line width=0.5pt] (1.00000,0.00000) -- (0.93795,0.99807);
\draw[black!35,line width=0.4pt] (0.50000,0.86603) -- (1.50000,0.86603);
\draw[red!70!black,line width=0.5pt] (0.50000,0.86603) -- (-0.45743,0.57735);
\draw[red!70!black,line width=0.5pt] (0.50000,0.86603) -- (1.33333,1.41880);
\draw[black!35,line width=0.4pt] (1.50000,0.86603) -- (0.77129,1.55084);
\draw[black!35,line width=0.4pt] (1.50000,0.86603) -- (1.56205,-0.13205);
\draw[black!35,line width=0.4pt] (1.50000,0.86603) -- (0.66667,0.31325);
\draw[red!70!black,line width=0.5pt] (1.50000,0.86603) -- (0.54257,0.57735);
\draw[black!35,line width=0.4pt] (0.83333,0.55277) -- (-0.06205,0.99807);
\draw[black!35,line width=0.4pt] (0.83333,0.55277) -- (0.77129,1.55084);
\draw[black!35,line width=0.4pt] (0.83333,0.55277) -- (1.56205,-0.13205);
\draw[black!35,line width=0.4pt] (0.83333,0.55277) -- (0.60462,-0.42072);
\draw[black!35,line width=0.4pt] (0.83333,0.55277) -- (-0.12409,0.26410);
\draw[black!35,line width=0.4pt] (0.83333,0.55277) -- (0.33333,1.41880);
\draw[red!70!black,line width=0.5pt] (0.83333,0.55277) -- (1.33333,1.41880);
\draw[black!35,line width=0.4pt] (-0.06205,0.99807) -- (0.77129,1.55084);
\draw[black!35,line width=0.4pt] (-0.06205,0.99807) -- (0.66667,0.31325);
\draw[black!35,line width=0.4pt] (-0.06205,0.99807) -- (-0.29076,0.02458);
\draw[red!70!black,line width=0.5pt] (-0.06205,0.99807) -- (0.43795,0.13205);
\draw[red!70!black,line width=0.5pt] (-0.06205,0.99807) -- (0.93795,0.99807);
\draw[black!35,line width=0.4pt] (0.77129,1.55084) -- (0.27129,0.68482);
\draw[red!70!black,line width=0.5pt] (0.77129,1.55084) -- (0.54257,0.57735);
\draw[black!35,line width=0.4pt] (1.56205,-0.13205) -- (0.72871,-0.68482);
\draw[black!35,line width=0.4pt] (1.56205,-0.13205) -- (0.60462,-0.42072);
\draw[black!35,line width=0.4pt] (1.56205,-0.13205) -- (0.66667,0.31325);
\draw[black!35,line width=0.4pt] (1.56205,-0.13205) -- (1.06205,0.73398);
\draw[black!35,line width=0.4pt] (0.72871,-0.68482) -- (0.66667,0.31325);
\draw[black!35,line width=0.4pt] (0.60462,-0.42072) -- (0.10462,0.44530);
\draw[black!35,line width=0.4pt] (0.60462,-0.42072) -- (-0.12409,0.26410);
\draw[black!35,line width=0.4pt] (0.60462,-0.42072) -- (-0.39538,-0.42072);
\draw[black!35,line width=0.4pt] (0.60462,-0.42072) -- (-0.29076,0.02458);
\draw[black!35,line width=0.4pt] (0.60462,-0.42072) -- (1.10462,0.44530);
\draw[red!70!black,line width=0.5pt] (0.60462,-0.42072) -- (0.54257,0.57735);
\draw[black!35,line width=0.4pt] (0.66667,0.31325) -- (-0.29076,0.02458);
\draw[black!35,line width=0.4pt] (0.66667,0.31325) -- (1.16667,-0.55277);
\draw[red!70!black,line width=0.5pt] (0.66667,0.31325) -- (0.16667,-0.55277);
\draw[black!35,line width=0.4pt] (0.10462,0.44530) -- (1.06205,0.73398);
\draw[black!35,line width=0.4pt] (0.10462,0.44530) -- (-0.39538,-0.42072);
\draw[black!35,line width=0.4pt] (0.10462,0.44530) -- (1.10462,0.44530);
\draw[black!35,line width=0.4pt] (0.10462,0.44530) -- (0.33333,1.41880);
\draw[red!70!black,line width=0.5pt] (0.10462,0.44530) -- (0.16667,-0.55277);
\draw[red!70!black,line width=0.5pt] (0.10462,0.44530) -- (-0.62409,1.13012);
\draw[red!70!black,line width=0.5pt] (0.10462,0.44530) -- (0.93795,0.99807);
\draw[black!35,line width=0.4pt] (-0.12409,0.26410) -- (0.37591,1.13012);
\draw[red!70!black,line width=0.5pt] (-0.12409,0.26410) -- (-0.62409,1.13012);
\draw[black!35,line width=0.4pt] (1.06205,0.73398) -- (0.33333,1.41880);
\draw[red!70!black,line width=0.5pt] (-0.39538,-0.42072) -- (-0.45743,0.57735);
\draw[red!70!black,line width=0.5pt] (-0.39538,-0.42072) -- (0.43795,0.13205);
\draw[red!70!black,line width=0.5pt] (-0.29076,0.02458) -- (0.54257,0.57735);
\draw[black!35,line width=0.4pt] (1.10462,0.44530) -- (1.16667,-0.55277);
\draw[black!35,line width=0.4pt] (1.10462,0.44530) -- (0.37591,1.13012);
\draw[black!35,line width=0.4pt] (1.10462,0.44530) -- (0.14719,0.15663);
\draw[red!70!black,line width=0.5pt] (1.10462,0.44530) -- (1.33333,1.41880);
\draw[red!70!black,line width=0.5pt] (0.33333,1.41880) -- (1.33333,1.41880);
\draw[red!70!black,line width=0.5pt] (0.33333,1.41880) -- (-0.62409,1.13012);
\draw[black!35,line width=0.4pt] (0.27129,0.68482) -- (0.04257,-0.28868);
\draw[red!70!black,line width=0.5pt] (0.27129,0.68482) -- (-0.62409,1.13012);
\draw[red!70!black,line width=0.5pt] (0.04257,-0.28868) -- (-0.45743,0.57735);
\draw[red!70!black,line width=0.5pt] (0.04257,-0.28868) -- (0.54257,0.57735);
\draw[red!70!black,line width=0.5pt] (1.16667,-0.55277) -- (0.43795,0.13205);
\draw[red!70!black,line width=0.5pt] (1.16667,-0.55277) -- (0.16667,-0.55277);
\draw[black!35,line width=0.4pt] (0.37591,1.13012) -- (0.14719,0.15663);
\draw[red!70!black,line width=0.5pt] (0.37591,1.13012) -- (-0.45743,0.57735);
\draw[red!70!black,line width=0.5pt] (0.37591,1.13012) -- (0.43795,0.13205);
\draw[red!70!black,line width=0.5pt] (0.37591,1.13012) -- (1.33333,1.41880);
\draw[red!70!black,line width=0.5pt] (0.37591,1.13012) -- (-0.62409,1.13012);
\draw[red!70!black,line width=0.5pt] (-0.45743,0.57735) -- (0.43795,0.13205);
\draw[red!70!black,line width=0.5pt] (-0.45743,0.57735) -- (0.54257,0.57735);
\draw[red!70!black,line width=0.5pt] (0.43795,0.13205) -- (0.93795,0.99807);
\fill[black] (0.00000,0.00000) circle (0.022);
\node[black,inner sep=1pt] at (-0.05689,-0.04887) {1};
\fill[black] (1.00000,0.00000) circle (0.022);
\node[black,inner sep=1pt] at (1.05925,-0.04598) {2};
\fill[black] (0.50000,0.86603) circle (0.022);
\node[black,inner sep=1pt] at (0.50414,0.94091) {3};
\fill[black] (1.50000,0.86603) circle (0.022);
\node[black,inner sep=1pt] at (1.56847,0.89663) {4};
\fill[black] (0.83333,0.55277) circle (0.022);
\node[black,inner sep=1pt] at (0.90286,0.58089) {5};
\fill[black] (-0.06205,0.99807) circle (0.022);
\node[black,inner sep=1pt] at (-0.11250,1.05357) {6};
\fill[black] (0.77129,1.55084) circle (0.022);
\node[black,inner sep=1pt] at (0.79012,1.62344) {7};
\fill[black] (1.56205,-0.13205) circle (0.022);
\node[black,inner sep=1pt] at (1.62923,-0.16539) {8};
\fill[black] (0.72871,-0.68482) circle (0.022);
\node[black,inner sep=1pt] at (0.74570,-0.75787) {9};
\fill[black] (0.60462,-0.42072) circle (0.022);
\node[black,inner sep=1pt] at (0.61624,-0.49482) {10};
\fill[black] (0.66667,0.31325) circle (0.022);
\node[black,inner sep=1pt] at (0.73396,0.28014) {11};
\fill[black] (0.10462,0.44530) circle (0.022);
\node[black,inner sep=1pt] at (0.03000,0.45288) {12};
\fill[black] (-0.12409,0.26410) circle (0.022);
\node[black,inner sep=1pt] at (-0.19702,0.24660) {13};
\fill[black] (1.06205,0.73398) circle (0.022);
\node[black,inner sep=1pt] at (1.12761,0.77039) {14};
\fill[black] (-0.39538,-0.42072) circle (0.022);
\node[black,inner sep=1pt] at (-0.44970,-0.47244) {15};
\fill[black] (-0.29076,0.02458) circle (0.022);
\node[black,inner sep=1pt] at (-0.35783,-0.00899) {16};
\fill[black] (1.10462,0.44530) circle (0.022);
\node[black,inner sep=1pt] at (1.17949,0.44977) {17};
\fill[black] (0.33333,1.41880) circle (0.022);
\node[black,inner sep=1pt] at (0.32295,1.49307) {18};
\fill[black] (0.27129,0.68482) circle (0.022);
\node[black,inner sep=1pt] at (0.22691,0.74528) {19};
\fill[black] (0.04257,-0.28868) circle (0.022);
\node[black,inner sep=1pt] at (0.00303,-0.35241) {20};
\fill[black] (1.16667,-0.55277) circle (0.022);
\node[black,inner sep=1pt] at (1.21051,-0.61362) {21};
\fill[black] (0.37591,1.13012) circle (0.022);
\node[black,inner sep=1pt] at (0.36575,1.20443) {22};
\fill[black] (0.14719,0.15663) circle (0.022);
\node[black,inner sep=1pt] at (0.08768,0.11099) {23};
\fill[red!80!black] (-0.45743,0.57735) circle (0.022);
\node[red!80!black,inner sep=1pt] at (-0.53122,0.59078) {24};
\fill[red!80!black] (0.43795,0.13205) circle (0.022);
\node[red!80!black,inner sep=1pt] at (0.42806,0.05770) {25};
\fill[red!80!black] (0.16667,-0.55277) circle (0.022);
\node[red!80!black,inner sep=1pt] at (0.14377,-0.62419) {26};
\fill[red!80!black] (0.54257,0.57735) circle (0.022);
\node[red!80!black,inner sep=1pt] at (0.57045,0.64698) {27};
\fill[red!80!black] (1.33333,1.41880) circle (0.022);
\node[red!80!black,inner sep=1pt] at (1.38188,1.47596) {28};
\fill[red!80!black] (-0.62409,1.13012) circle (0.022);
\node[red!80!black,inner sep=1pt] at (-0.68676,1.17132) {29};
\fill[red!80!black] (0.93795,0.99807) circle (0.022);
\node[red!80!black,inner sep=1pt] at (0.98426,1.05707) {30};
\end{tikzpicture}

\caption{The configuration $X_{\NPTS}$ and its $\NUNIT$ unit distances. The $23$ points of \cite{ACMVZ} are black, the
$\NNEW$ new points red; red segments are the unit distances at a new point.}\label{fig:config}
\end{figure}

\section{The certificate and its verification}\label{sec:checks}

The program for $X_{\NPTS}$ has the inclusion--exclusion rows of \cite{ACMVZ}: one variable for each independent set,
the equations $\Pr(i\in S)=\delta$ and $\Pr(i,j\in S)=f(r_{ij})$, and the congruence relations. We generate the latter
ourselves: the independent sets with at least three points fall into congruence classes (compared exactly), and each
class $S_1,\dots,S_m$ gives the $m-1$ relations $\Pr(S_1\subseteq S)=\Pr(S_j\subseteq S)$, with $S_j$ listed in the order
of a distance-preserving bijection. (Relations between pairs repeat the equations $\Pr(i,j\in S)=f(r_{ij})$.)

The certificate is found in two steps. First the program is solved in floating point, with $\kappa$ supported on the
grid $t\in\{0,0.1,\dots,100\}\cup\{100.25,100.5,\dots,400\}$; its multipliers $w$ of the relations are rounded to
multiples of $2^{-40}$ and frozen.
For frozen $w$ the best $\psi_0,\psi_T,\lambda,u$ solve a second, small program (one variable per point and per pair):
minimize $\psi_T$ subject to $Q\ge0$ on every independent set, $-\sum_i\lambda_i\ge1.03$, and
$\Phi(t)\ge1+2\cdot10^{-5}$ at the points of a grid of step $0.005$ on $[0,50]$ and $0.02$ on $[50,1000]$ where it is
needed (rows are added at the local minima of $\Phi$ until none is violated). Then $\psi_0$, $\lambda$ and $u$ are
rounded to multiples of $2^{-40}$ and $\psi_T$ is raised until $Q\ge0$ holds exactly on every independent set. The
certificate for $X_{\NPTS}$ uses $\NCONG$ relations and gives $\psi_T<\MONE$; it is the file
\texttt{data/forbidden\_distances/m1/certificates/m1\_x30.json}. Nothing in this construction needs to be trusted: the checks below take the certificate as it is.

Two programs, written separately, check the certificate and prove Theorem~\ref{thm:main}:
\begin{itemize}
\item \texttt{check\_m1.py} computes the unit-distance graph of $X$ exactly in $\Q(\sqrt3,\sqrt{11})$, checks that every
index lies in range, that no $u_{ij}$ sits on a unit pair and that every relation in $\mathcal C$ is a congruence, checks
$Q\ge0$ on all independent sets with integers over a common denominator, and proves $\Phi\ge1$ on $[0,T]$ by bisection
with Taylor enclosures of $J_0$ in Arb ball arithmetic, and beyond $T$ by $|J_0(x)|\le\sqrt{2/(\pi x)}$; it
cross-checks the number of independent sets by the recursion $I(G)=I(G-v)+I(G-N[v])$, and it does not rely on
\texttt{assert} statements, which Python skips under its option \texttt{-O}.
\item \texttt{check\_m1\_indep.py} has its own exact arithmetic, its own transcription-independent reading of the points,
its own integer interval arithmetic with the power series and the Hankel expansion of $J_0$, and an adaptive sweep based
on $|J_0''|\le1$; it computes $Q$ on all independent sets by a subset-sum pass and recomputes a sample from the
definition.
\end{itemize}
The facts about $J_0$ used are: $J_0(x)=\frac1\pi\int_0^\pi\cos(x\sin s)\,ds$, so $|J_0|\le1$ and $|J_0''|\le1$;
$J_0\ge-0.4028$, proved by \texttt{check\_m1.py} on $[0,4]$ by interval arithmetic and beyond by the next bound; and
$|J_0(x)|\le\sqrt{2/(\pi x)}$ for $x>0$, which follows from Nicholson's formula
$J_0(x)^2+Y_0(x)^2=\frac8{\pi^2}\int_0^\infty K_0(2x\sinh u)\,du$ \cite[10.9.30]{DLMF}, \cite[§§13.73--13.74]{Watson}.

Both accept the certificate (in under a minute) and the certificate for $X_{23}$ alone, which gives $0.24698967$ and
reproduces the bound of \cite{ACMVZ}. Both reject corrupted copies: an index out of range (also in one-point
relations), a broken congruence, a $u_{ij}$ on a unit pair, a moved point, a lowered $\psi_T$, a raised $\lambda_i$, a
lowered $\psi_0$. A referee report, written by a separate AI agent, proves Proposition~\ref{prop:ie} in full; its own checker, written before it read ours, accepts the four certificates, and of $28$ corrupted certificates every invalid one is rejected by all three checkers. It found one weakness, now removed: \texttt{check\_m1.py} relied on \texttt{assert} statements.

\section{How the points were found}\label{sec:search}

The search follows \cite[§7]{ACMVZ}: a new point must lie at distance $1$ from at least two points already present and at
distance at least $0.1$ from all of them. For $X_{23}$ there are $208$ such points. Adding one of them alone lowers the
value of the program (in floating point, discretized with step $0.1$ up to $t=100$, where $X_{23}$ gives $0.246948$) to as
little as $0.246083$, and the gains of different points partly add up: the two best together gain $1.36\cdot10^{-3}$,
four fifths of the sum of their single gains. The six best points at distance $1$ from at least four points of
$X_{23}$, added in the order of their single gains, give $X_{29}$, and $x_{30}$ is the best of the others, with the third
largest single gain (Table~\ref{tab:values}). Each point gains about $4\cdot10^{-4}$, slowly less as the set grows, while the program grows with
the number of independent sets; we stopped at $\NPTS$ points for the certificate.

\begin{table}[ht]
\centering
\begin{tabular}{rrrcc}
\toprule
points & independent sets & relations & value (floating point) & certified bound\\
\midrule
$23$ & 13\,552 & 2\,122 & $0.246948$ & $0.24698967$ \\
$25$ & 19\,868 & 4\,091 & $0.245583$ & -- \\
$26$ & 28\,622 & 5\,008 & $0.245169$ & $0.24522384$ \\
$27$ & 33\,898 & 6\,044 & $0.244841$ & -- \\
$29$ & 55\,610 & 10\,012 & $0.244092$ & $0.24414455$ \\
$30$ & 84\,893 & 15\,709 & $0.243743$ & $0.24379682$ \\
\bottomrule
\end{tabular}
\caption{The program for $X_{23}$ and its extensions: the value in floating point (with $\kappa$ on the grid of Section~\ref{sec:checks}), and the bound proved by a certificate (rounded up).}\label{tab:values}
\end{table}

\begin{remark}
Ambrus et al.\ observed that the last point of their $24$-point set did not improve the bound and removed it
\cite[§9]{ACMVZ}. Our points improve it from the first one on: the best single point lowers the value by $8.7\cdot10^{-4}$.
We do not know how the announced bound $0.2415$ \cite{Varga} was obtained, so we cannot compare the methods.
\end{remark}

\begin{thebibliography}{99}
\bibitem{ACMVZ} G. Ambrus, A. Csisz\'arik, M. Matolcsi, D. Varga and P. Zs\'amboki, \emph{The density of planar sets
avoiding unit distances}, Math. Program. 207 (2024) 303--327; arXiv:2207.14179.
\bibitem{AM} G. Ambrus and M. Matolcsi, \emph{Density estimates of $1$-avoiding sets via higher order correlations},
Discrete Comput. Geom. 67 (2022) 1245--1256; arXiv:1809.05453.
\bibitem{BPS} T. Bellitto, A. P\^echer and A. S\'edillot, \emph{On the density of sets of the Euclidean plane avoiding
distance $1$}, Discrete Math. Theor. Comput. Sci. 23(1) (2021); arXiv:1810.00960.
\bibitem{Croft} H. T. Croft, \emph{Incidence incidents}, Eureka 30 (1967) 22--26.
\bibitem{DLMF} \emph{NIST Digital Library of Mathematical Functions}, \url{https://dlmf.nist.gov}.
\bibitem{DV} \'A. D\'ucz and D. Varga, \emph{A unit-distance graph in the plane with independence ratio below $1/4$},
arXiv:2606.28157 (26 June 2026).
\bibitem{Erdos85} P. Erd\H{o}s, \emph{Problems and results in combinatorial geometry}, in: Discrete Geometry and
Convexity, Ann. New York Acad. Sci. 440 (1985) 1--11.
\bibitem{Galan} S. Gal\'an, \emph{Measurable density bounds hold for all colourings of the plane with algebraic
forbidden distances}, preprint, version 8 (11 October 2026),
\url{https://github.com/decalion89/chromatic-number-of-the-plane}.
\bibitem{KMOR} T. Keleti, M. Matolcsi, F. M. de Oliveira Filho and I. Z. Ruzsa, \emph{Better bounds for planar sets
avoiding unit distances}, Discrete Comput. Geom. 55 (2016) 642--661; arXiv:1501.00168.
\bibitem{OFV} F. M. de Oliveira Filho and F. Vallentin, \emph{Fourier analysis, linear programming, and densities of
distance avoiding sets in $\R^n$}, J. Eur. Math. Soc. 12 (2010) 1417--1428.
\bibitem{OpenAI} OpenAI, \emph{The Euclidean plane is not five-colorable}, preprint (23 September 2026), with a
formalization in Lean, \url{https://github.com/openai/math}.
\bibitem{Varga} D. Varga, \emph{Solving a conjecture of Erd\H{o}s}, blog post of the AI group of the Alfr\'ed R\'enyi
Institute of Mathematics, 10 December 2023, \url{https://ai.renyi.hu/posts/solving-erdos-conjecture/}.
\bibitem{Watson} G. N. Watson, \emph{A Treatise on the Theory of Bessel Functions}, 2nd ed., Cambridge University
Press, 1944.
\bibitem{Szekely} L. A. Sz\'ekely, \emph{Measurable chromatic number of geometric graphs and sets without some distances
in Euclidean space}, Combinatorica 4 (1984) 213--218.
\end{thebibliography}

\end{document}
